\answer{OP25}
任意两条弦的四个端点几乎必然互异。条件于四个端点的圆周顺序，它们形成两条弦的三种无序配对等可能。只有连接交错端点的那一种配对在圆内部相交，因此每对弦相交的概率为 $1/3$。

记 $I_{ij}$ 为第 $i,j$ 条弦在圆内部相交的指标。连续取样下，端点重合、弦重合以及三条弦共点均是零概率事件；例如先固定两条相交弦，再取第三条弦的一个端点，其另一个端点只有至多一个位置能使它通过既有交点。故交点数几乎必然等于
\[
K=\sum_{1\le i<j\le n}I_{ij}.
\]
无需假设这些指标独立，直接得到
\[
\boxed{\E K=\binom n2\frac13=\frac{n(n-1)}6.}
\]
此处的随机弦模型由两个独立均匀端点定义；换一种选弦方式可能改变答案。
